\documentclass[11pt,a4paper]{article}

\usepackage[utf8]{inputenc}
\usepackage[T1]{fontenc}
\usepackage{lmodern}
\usepackage{amsmath,amssymb,amsthm}
\usepackage{booktabs}
\usepackage{enumitem}
\usepackage{hyperref}
\usepackage[margin=2.5cm]{geometry}
\usepackage{xcolor}

\definecolor{codeblue}{RGB}{30,90,156}

\newtheorem{theorem}{Theorem}[section]
\newtheorem{proposition}[theorem]{Proposition}
\newtheorem{corollary}[theorem]{Corollary}
\newtheorem{lemma}[theorem]{Lemma}

\theoremstyle{definition}
\newtheorem{definition}[theorem]{Definition}

\theoremstyle{remark}
\newtheorem*{remark}{Remark}

\hypersetup{
  colorlinks=true,
  linkcolor=codeblue,
  citecolor=codeblue,
  urlcolor=codeblue,
}

\newcommand{\Nd}{\mathbb{N}_\delta}
\newcommand{\Zd}{\mathbb{Z}_\delta}
\newcommand{\Qd}{\mathbb{Q}_\delta}
\newcommand{\Sd}{S_\delta}
\newcommand{\zd}{0_\delta}
\newcommand{\oad}{\oplus}
\newcommand{\omd}{\otimes}

\title{%
  Arithmetic from Distinction:\\
  A Number Tower Forced by a Single Primitive}
\author{%
  Jonathan Washburn\\
  Recognition Physics Institute, Austin, TX, USA\\
  \texttt{jon@recognitionphysics.org}}
\date{May 2026}

\begin{document}

\maketitle

\begin{abstract}\sloppy
We construct the natural numbers, the integers, and the rational field from a
single primitive: the act of distinction. The construction uses no axiom of
infinity in the set-theoretic sense, no axiom of choice, and no law of excluded
middle. Each object is a finite record of performed distinctions, and every
theorem is proved by structural induction over such finite records. From the
primitive we build finite traces, a sameness relation that certifies the
absence of a separating distinction, a difference relation that certifies its
presence, and a quotient discipline under which any admissible predicate must
respect sameness. Counting is the length of a distinction orbit. Sign is the
balance of two orbit lengths. Ratio is the cross-comparison of two signed
orbits against two lengths. On these objects we prove the full slate of
elementary arithmetic: associativity, commutativity, distributivity,
cancellation, a decidable total order, absolute value, divisibility, primality,
Euclidean division with unique quotient and remainder, greatest common divisors,
coprimality, and the field axioms for the rationals. Each construction carries
a proof-strength tag, and every result in this paper carries the strongest tag
available: it is forced by distinction and finite repetition alone. To certify
that the construction recovers the familiar objects and nothing stranger, we
exhibit faithful display maps onto the standard naturals, integers, and
rationals, and we prove these maps are isomorphisms of the relevant algebraic
structures. The display maps are verification instruments, not ingredients:
the orbit operations are defined without reference to any prior numeral system.
The upshot is that the elementary number tower is not a convention chosen on
top of a foundation but a structure that the single act of distinction permits
and then forces. We close by marking precisely where this forcing stops. The
completion of the rational field into the continuum is a separate commitment,
treated in a companion paper, and it is not a consequence of distinction.
\end{abstract}

\tableofcontents

%% ===================================================================
\section{What this paper does}\label{sec:intro}

Foundations of mathematics usually begin with a rich primitive and then encode
number inside it. Set theory begins with membership and builds the number three
as a particular nest of braces. Type theory begins with types and terms and
introduces the natural numbers as a particular inductive type. Both encodings
work, and both are correct. Both also carry a feature that is easy to overlook
because it is so familiar: the primitive has nothing to do with counting. A set
is not a number. A type is not a number. The number arrives only after a chain
of definitions that translate the foreign primitive into arithmetic behavior.

This paper takes a different primitive. We begin with the act of distinction,
the operation of marking that one thing is not another, and we show that the
elementary number tower follows from it directly. Counting is the most natural
reading of repeated distinction: to count is to perform the same act of marking
again and again and to keep a record of how many times. Sign is the most
natural reading of two opposed directions of counting. Ratio is the most natural
reading of comparing two counts. None of these is imposed as a definitional
convention. Each is the structure that the primitive already contains once you
take it seriously.

The reason to care is not that we obtain a new arithmetic. The integers built
here are isomorphic to the usual integers, and the rationals to the usual
rationals; we prove exactly that. The reason to care is the audit. Every result
in this paper is established with a single, fixed, minimal proof budget, and we
keep an explicit ledger of that budget. We use no axiom of infinity in the
set-theoretic sense, meaning we never posit a completed infinite set; we work
only with finite records and the two rules that generate them. We use no axiom
of choice, because every object we need is constructed explicitly rather than
selected from a family. We use no law of excluded middle, because every
predicate we rely on is decidable by a terminating procedure on finite data. A
foundation that announces a single primitive owes the reader a precise account
of what that primitive can and cannot reach without help. This paper is that
account for the arithmetic layer.

\subsection{Contributions}

The paper makes four contributions.

\begin{enumerate}[leftmargin=2em]
  \item It defines a finite kernel from the primitive act of distinction:
    traces, trace concatenation with its associativity law, the sameness and
    difference relations with their consistency condition, and the quotient
    discipline that any admissible predicate must obey.
  \item It constructs the natural numbers as distinction orbits and proves the
    full arithmetic slate for them, including a decidable total order and a
    cancellation law, by structural induction alone.
  \item It constructs the integers as balanced pairs of orbit lengths and the
    rationals as cross-compared ratios, and it proves that these carry a
    commutative ring structure and an ordered field structure respectively,
    each isomorphic to the standard object under a faithful display map.
  \item It derives the elementary number-theoretic surface, divisibility,
    primality, Euclidean division, greatest common divisors, and coprimality,
    with the same minimal proof budget, and it states exactly where the budget
    would have to grow, namely at the completion into the continuum, which is
    deferred to a companion paper.
\end{enumerate}

\subsection{Status discipline}

A paper that claims to derive arithmetic from one act, rather than to encode it
inside a richer theory, must be careful about what counts as derivation and what
smuggles in a hidden assumption. We adopt a per-theorem discipline. Every
construction is labeled with a proof-strength tag, defined in
Section~\ref{sec:strength}, and every result in this paper carries the same tag:
it is forced by distinction and finite repetition alone, with constructive logic
and no further commitments. When we use a display map into a standard object, we
say so explicitly and we explain that the map is used to certify faithfulness
after the fact, not to define the construction. The reader should be able to
check at every step that nothing has entered the argument except the primitive
and induction over finite records.

%% ===================================================================
\section{The primitive: distinction acts and traces}\label{sec:kernel}

The primitive is an act, not an object. It is the operation of marking a
difference, taken in its most contentless form.

\begin{definition}[Distinction act]\label{def:act}
A \emph{distinction act} is a single, atomic operation of marking a difference.
It carries no parameters, no payload, and no internal structure. It is pure
separation: the minimal operation that takes an undifferentiated situation to a
differentiated one. We write $\delta$ for this act.
\end{definition}

A distinction act on its own is too thin to support mathematics. What supports
mathematics is the record of having performed distinction acts, possibly many
times. That record is a trace.

\begin{definition}[Trace]\label{def:trace}
A \emph{trace} is a finite sequence of distinction acts, defined by two
formation rules:
\begin{enumerate}[leftmargin=2em, label=\textup{(\roman*)}]
  \item the \emph{empty trace} $\varepsilon$ is a trace, recording that no
    distinction has been performed;
  \item if $\tau$ is a trace, then $\tau \cdot \delta$, the result of extending
    $\tau$ by one further distinction act, is a trace.
\end{enumerate}
The \emph{length} of a trace is the number of extension steps used to build it:
$|\varepsilon| = 0$ and $|\tau \cdot \delta| = |\tau| + 1$. The
\emph{concatenation} $\tau_1 \mathbin{\|} \tau_2$ of two traces is defined by
recursion on the second argument:
\[
  \tau_1 \mathbin{\|} \varepsilon = \tau_1,
  \qquad
  \tau_1 \mathbin{\|} (\tau_2 \cdot \delta) = (\tau_1 \mathbin{\|} \tau_2) \cdot \delta .
\]
\end{definition}

Two features of this definition deserve emphasis, because they carry the whole
budget claim of the paper. First, every trace is finite. There is no rule that
produces an infinite trace, and the two rules above generate only objects built
from finitely many extension steps. When we later quantify over all traces, we
quantify over a collection generated by these rules, exactly as one quantifies
over the natural numbers generated by zero and successor. We never posit a
completed infinite totality as a new object with its own axiom. Second, the
primitive is an act, not a set, a type, a proposition, or a field. Set theory
tells you what a set is; type theory tells you what a type is; the present
theory tells you what distinguishing is, and lets sets, types, and numbers
appear later as constructions rather than as starting assumptions.

\begin{lemma}[Concatenation is associative]\label{lem:concat}
For all traces $\tau_1, \tau_2, \tau_3$,
\[
  (\tau_1 \mathbin{\|} \tau_2) \mathbin{\|} \tau_3
  = \tau_1 \mathbin{\|} (\tau_2 \mathbin{\|} \tau_3).
\]
\end{lemma}

\begin{proof}
By structural induction on $\tau_3$. If $\tau_3 = \varepsilon$, both sides equal
$\tau_1 \mathbin{\|} \tau_2$, using the base clause of concatenation twice. For
the inductive step, suppose the identity holds for $\tau_3$ and consider
$\tau_3 \cdot \delta$. Using the recursive clause of concatenation on each side
and the induction hypothesis,
\begin{align*}
  (\tau_1 \mathbin{\|} \tau_2) \mathbin{\|} (\tau_3 \cdot \delta)
   &= \bigl( (\tau_1 \mathbin{\|} \tau_2) \mathbin{\|} \tau_3 \bigr) \cdot \delta
    = \bigl( \tau_1 \mathbin{\|} (\tau_2 \mathbin{\|} \tau_3) \bigr) \cdot \delta \\
   &= \tau_1 \mathbin{\|} \bigl( (\tau_2 \mathbin{\|} \tau_3) \cdot \delta \bigr)
    = \tau_1 \mathbin{\|} \bigl( \tau_2 \mathbin{\|} (\tau_3 \cdot \delta) \bigr).
\end{align*}
This is the claim for $\tau_3 \cdot \delta$, completing the induction.
\end{proof}

%% ===================================================================
\section{Sameness, difference, and the quotient discipline}\label{sec:samediff}

Traces let us record distinctions, but mathematics also needs the converse: a
way to certify that two things have not been distinguished. In set theory,
equality is a logical primitive fixed by extensionality. In type theory, it is
the identity type. In a theory built on acts of distinction, the natural
primitive is a certification relative to a trace: two endpoints are the same
when no distinction in the current trace separates them, and different when one
does.

\begin{definition}[Sameness and difference]\label{def:samediff}
Let $A$ be a collection of trace endpoints. We introduce two relations on $A$.
The \emph{sameness} relation $a \sim b$ asserts that no distinction in the
current trace separates $a$ from $b$. The \emph{difference} relation
$a \mathbin{\#} b$ asserts that some distinction in the trace does separate them.
We require $\sim$ to be an equivalence relation: reflexive, symmetric, and
transitive. We require the \emph{consistency condition}, that for any pair
$(a,b)$ and any fixed respect, sameness and difference cannot both hold,
\[
  a \sim b \;\wedge\; a \mathbin{\#} b \;\Longrightarrow\; \bot .
\]
\end{definition}

The relation $\sim$ is weaker than logical equality, and this is the point. Two
endpoints can be the same relative to a coarse trace and different relative to a
finer one. Sameness is a statement about what the present record can resolve,
not a metaphysical claim about identity. This is the right primitive for a
theory whose only generator is the act of distinction: you are entitled to treat
two things as the same exactly when you have performed no act that tells them
apart.

Once endpoints are identified under sameness, every construction built on top
must respect that identification. Otherwise the construction would distinguish
what the trace cannot, which violates the spirit of the primitive.

\begin{definition}[Admissible predicate and quotient descent]\label{def:subst}
A predicate $P$ on endpoints is \emph{admissible} if it respects sameness:
whenever $a \sim b$, we have $P(a)$ if and only if $P(b)$. An admissible
predicate descends to a well-defined predicate on the quotient
$A / {\sim}$. More generally, a function $f$ on representatives descends to a
function on $A / {\sim}$ whenever $a \sim b$ implies $f(a)$ and $f(b)$ agree in
the relevant sense. This descent principle is the only tool we use to move from
representatives to quotient objects, and it is available because $\sim$ is an
equivalence relation.
\end{definition}

The kernel is now complete. It consists of finite traces with their
concatenation and associativity law, the sameness and difference relations with
their consistency condition, and the descent principle for admissible predicates
and functions. Everything in the rest of the paper is built from this kernel by
finite construction and structural induction.

%% ===================================================================
\section{The proof-strength ledger}\label{sec:strength}

Conventional mathematics answers the question ``what are your assumptions?''
once, globally, and usually with a large answer such as set theory together with
choice and classical logic. We answer it per theorem, with a tag attached to
each result. The tags form a lattice ordered by how much machinery each level
adds beyond the bare kernel.

\begin{table}[ht]
\centering
\begin{tabular}{ll}
\toprule
\textbf{Tag} & \textbf{Meaning} \\
\midrule
distinction-only &
  Uses only the finite kernel and constructive logic. No infinity, choice, or excluded middle. \\
distinction $+$ closure &
  Adds completed infinite streams or limiting boundary commitments. \\
distinction $+$ choice &
  Uses a choice principle to select representatives or witnesses. \\
classical extension &
  Uses classical logic or an imported classical carrier such as the real line. \\
\bottomrule
\end{tabular}
\caption{The proof-strength ledger. Each level includes the machinery of all
levels below it. Every theorem in this paper carries the lowest tag,
distinction-only.}
\label{tab:strength}
\end{table}

The discipline costs the author bookkeeping and gives the reader something most
foundations do not provide: granular proof-strength information at the level of
the individual theorem. A reader who cares only about constructive content can
read the distinction-only results and stop. A reader who wants to know exactly
where a completed infinity or a choice principle would first be required can
find that boundary marked. In this paper the boundary is never crossed: the
arithmetic tower through the rational field lives entirely at the
distinction-only level. The first genuine crossing, into the continuum, is the
subject of the companion paper and is not attempted here.

%% ===================================================================
\section{Base-neutral number}\label{sec:number}

Number is the length of a record of repeated distinction. We call such a record
a distinction orbit, and we treat the orbit itself as the number, prior to any
numeral system used to display it.

\begin{definition}[Distinction orbit, $\Nd$]\label{def:orbit}
The collection $\Nd$ of \emph{distinction orbits} is generated by two rules:
\begin{enumerate}[leftmargin=2em, label=\textup{(\roman*)}]
  \item the \emph{zero orbit} $\zd$ is an orbit, recording that no distinction
    has been performed;
  \item if $n$ is an orbit, then its \emph{successor} $\Sd(n)$, recording one
    further distinction, is an orbit.
\end{enumerate}
The \emph{display map} $\varphi \colon \Nd \to \mathbb{N}$ onto the standard
naturals is fixed by $\varphi(\zd) = 0$ and $\varphi(\Sd(n)) = \varphi(n) + 1$.
\end{definition}

The phrase \emph{base-neutral} records that the orbit precedes every numeral
system. The decimal symbol $3$, the binary symbol $11$, the Roman symbol
$\mathrm{III}$, and the successor expression $\Sd(\Sd(\Sd(\zd)))$ are four
display surfaces for one object: three performed distinctions. None of those
symbols is the number; each is a way of writing it. The orbit is the number, and
the display map converts it into whichever numeral system a given application
finds convenient. Tally marks on a wall are closer to the orbit than the
positional decimal system is, because tally marks are direct records of
performed acts and carry no positional weighting, no carry rule, and no
zero-padding. Those features belong to the numeral system, not to the number.

We define arithmetic by recursion on orbit structure, never referring to any
numeral system.

\begin{definition}[Orbit addition and multiplication]\label{def:orbitops}
Addition $\oad$ and multiplication $\omd$ on $\Nd$ are defined by recursion on
the second argument:
\begin{align}
  n \oad \zd &= n, &
  n \oad \Sd(m) &= \Sd(n \oad m); \label{eq:add}\\
  n \omd \zd &= \zd, &
  n \omd \Sd(m) &= (n \omd m) \oad n. \label{eq:mul}
\end{align}
\end{definition}

Addition is concatenation of orbits: appending $m$ further distinctions after
$n$. Multiplication is repeated addition: $n \omd m$ adds $n$ to itself $m$
times. Both are defined by the two formation rules alone and are therefore
distinction-only.

\begin{theorem}[Orbit arithmetic]\label{thm:orbit}
The structure $(\Nd, \zd, \Sd, \oad, \omd)$ satisfies, with distinction-only
strength:
\begin{enumerate}[leftmargin=2em, label=\textup{(\alph*)}]
  \item additive and multiplicative associativity,
    $(a \oad b) \oad c = a \oad (b \oad c)$ and
    $(a \omd b) \omd c = a \omd (b \omd c)$;
  \item additive and multiplicative commutativity,
    $a \oad b = b \oad a$ and $a \omd b = b \omd a$;
  \item distributivity, $a \omd (b \oad c) = (a \omd b) \oad (a \omd c)$;
  \item identities, $a \oad \zd = a$ and $a \omd \Sd(\zd) = a$;
  \item additive cancellation, $a \oad c = b \oad c \Rightarrow a = b$;
  \item display faithfulness, $\varphi(a \oad b) = \varphi(a) + \varphi(b)$ and
    $\varphi(a \omd b) = \varphi(a) \cdot \varphi(b)$.
\end{enumerate}
\end{theorem}

\begin{proof}
Each item is proved by structural induction on $\Nd$, using only the recursive
clauses~\eqref{eq:add} and~\eqref{eq:mul}. We give the additive arguments in
full and indicate the rest, since all follow the same pattern.

First, a base lemma: $\zd \oad n = n$ for all $n$. Induct on $n$. The case
$n = \zd$ is the base clause of addition. If $\zd \oad m = m$, then
$\zd \oad \Sd(m) = \Sd(\zd \oad m) = \Sd(m)$. Next, a sliding lemma:
$\Sd(a) \oad m = \Sd(a \oad m)$. Induct on $m$. For $m = \zd$, both sides are
$\Sd(a)$. If $\Sd(a) \oad m = \Sd(a \oad m)$, then
$\Sd(a) \oad \Sd(m) = \Sd(\Sd(a) \oad m) = \Sd(\Sd(a \oad m)) = \Sd(a \oad \Sd(m))$.

Additive associativity follows by induction on $c$. For $c = \zd$, both sides
are $a \oad b$. The inductive step uses the recursive clause three times and the
hypothesis once. Additive commutativity follows by induction on $b$, using the
base lemma for $b = \zd$ and the sliding lemma in the step:
$a \oad \Sd(m) = \Sd(a \oad m) = \Sd(m \oad a) = \Sd(m) \oad a$. Additive
cancellation follows by induction on $c$: the case $c = \zd$ is immediate, and
the step uses the injectivity of $\Sd$, which holds because $\Sd$ is a formation
rule with no identifications. Multiplicative associativity, multiplicative
commutativity, and distributivity follow by the same method, each reducing to
the additive laws already established. Display faithfulness is immediate from
the recursive clauses and the recursive definition of $\varphi$.
\end{proof}

We also fix the order on orbits, which we use throughout the sign and ratio
constructions.

\begin{definition}[Order and truncated subtraction on $\Nd$]\label{def:order}
The relation $\leq_\delta$ on $\Nd$ is generated by $\zd \leq_\delta n$ for all
$n$, and $\Sd(m) \leq_\delta \Sd(n)$ if and only if $m \leq_\delta n$. Truncated
subtraction $\ominus$ is defined by $n \ominus \zd = n$,
$\zd \ominus \Sd(m) = \zd$, and $\Sd(n) \ominus \Sd(m) = n \ominus m$.
\end{definition}

\begin{proposition}[The order is a decidable total order]\label{prop:order}
The relation $\leq_\delta$ is reflexive, antisymmetric, transitive, and total,
and for any two orbits $m, n$ the question whether $m \leq_\delta n$ is decided
by a terminating recursion. Moreover $\leq_\delta$ is compatible with addition:
$m \leq_\delta n$ if and only if $m \oad c \leq_\delta n \oad c$.
\end{proposition}

\begin{proof}
Each property is proved by induction following the recursive clauses of
Definition~\ref{def:order}. The comparison recursion strips one successor from
each argument at each step and halts when one argument reaches $\zd$, so it
terminates after a number of steps equal to the smaller length. Because the
recursion returns a definite yes or no for every input, the order is decidable
without any appeal to excluded middle. Compatibility with addition follows by
induction on $c$.
\end{proof}

The decidability claim in Proposition~\ref{prop:order} is the first concrete
payoff of the budget discipline. In a classical setting one would assert that
$m \leq_\delta n$ or $n <_\delta m$ by the law of excluded middle and move on.
Here the disjunction is not asserted; it is computed. The comparison is a finite
procedure on finite records, so the trichotomy holds constructively.

%% ===================================================================
\section{Sign as balance: the integers}\label{sec:sign}

The natural numbers have no sign. To subtract, one needs a direction. The usual
construction defines an integer as an equivalence class of pairs of naturals
under the relation $(a,b) \sim (c,d)$ exactly when $a + d = b + c$. The
construction is correct, but the reason an integer points in a direction is
invisible in it; one must know the convention that the first coordinate is the
positive part. In the present setting, sign is balance, and the direction is
visible in the object.

\begin{definition}[Signed orbit and balanced equality]\label{def:signedorbit}
A \emph{signed orbit} is an ordered pair $(p, n) \in \Nd \times \Nd$, where $p$
is the positive part and $n$ the negative part. Two signed orbits $(a_+, a_-)$
and $(b_+, b_-)$ are \emph{balanced-equal}, written
$(a_+, a_-) \approx (b_+, b_-)$, when
\[
  a_+ \oad b_- \;=\; b_+ \oad a_-
\]
holds in $\Nd$. The \emph{display} of a signed orbit is
$\psi(p, n) = \varphi(p) - \varphi(n) \in \mathbb{Z}$.
\end{definition}

The balance condition is the cross-addition condition, written entirely inside
orbit arithmetic. It compares the two pairs without ever subtracting and without
ever mentioning a negative number, which is exactly what one wants at the level
where negative numbers are being constructed rather than assumed. For example,
the integer $-1$ is any signed orbit $(p, n)$ whose negative part exceeds its
positive part by one performed distinction, such as $(\zd, \Sd(\zd))$ or
$(\Sd(\zd), \Sd(\Sd(\zd)))$, and balance identifies all of them.

\begin{definition}[The integers $\Zd$]\label{def:prcint}
Set $\Zd = (\Nd \times \Nd) / {\approx}$, the quotient of signed orbits by
balanced equality. Operations descend from representatives:
\begin{align}
  [(a_+, a_-)] + [(b_+, b_-)] &= [(a_+ \oad b_+,\; a_- \oad b_-)],
    \label{eq:intadd}\\
  -[(a_+, a_-)] &= [(a_-, a_+)], \label{eq:intneg}\\
  [(a_+, a_-)] \cdot [(b_+, b_-)] &=
    [(\,a_+ \omd b_+ \oad a_- \omd b_-,\;\; a_+ \omd b_- \oad a_- \omd b_+\,)].
    \label{eq:intmul}
\end{align}
\end{definition}

Negation swaps the positive and negative parts: it reverses direction, which is
exactly what a sign change should mean for a balance. Addition adds the two
positive parts and the two negative parts separately. Multiplication follows the
familiar sign rule, written so that the positive part of a product collects the
two like-signed cross terms and the negative part collects the two opposite-signed
cross terms.

\begin{theorem}[The integers form a ring isomorphic to $\mathbb{Z}$]\label{thm:int}
With distinction-only strength: the relation $\approx$ is an equivalence
relation; the operations \eqref{eq:intadd} through \eqref{eq:intmul} are
well-defined on the quotient; $\Zd$ with these operations is a commutative ring
with additive identity $[(\zd, \zd)]$ and multiplicative identity
$[(\Sd(\zd), \zd)]$; and the display map $\Psi \colon \Zd \to \mathbb{Z}$,
$\Psi([(p, n)]) = \varphi(p) - \varphi(n)$, is a ring isomorphism.
\end{theorem}

\begin{proof}
That $\approx$ is an equivalence relation follows from the additive laws of
Theorem~\ref{thm:orbit}. Reflexivity is $a_+ \oad a_- = a_+ \oad a_-$. Symmetry
is immediate. For transitivity, suppose $a_+ \oad b_- = b_+ \oad a_-$ and
$b_+ \oad c_- = c_+ \oad b_-$. Adding the two equalities and using
commutativity and associativity gives
$(a_+ \oad c_-) \oad (b_+ \oad b_-) = (c_+ \oad a_-) \oad (b_+ \oad b_-)$, and
additive cancellation, Theorem~\ref{thm:orbit}(e), removes the common term
$b_+ \oad b_-$ to yield $a_+ \oad c_- = c_+ \oad a_-$.

Well-definedness of the operations is the standard cross-addition computation.
For addition, if $(a_+, a_-) \approx (a_+', a_-')$ and
$(b_+, b_-) \approx (b_+', b_-')$, adding the two balance equalities and
rearranging gives the balance equality for the sums. For negation, swapping
parts turns one balance equality into the other. For multiplication, one expands
both sides using distributivity and the balance hypotheses; the computation is
finite and uses only Theorem~\ref{thm:orbit}. The ring axioms transfer from
$\Nd$ in the same way: each is verified on representatives and respects balance.

The display map $\Psi$ is well-defined because balanced orbits have equal
display: $a_+ \oad b_- = b_+ \oad a_-$ implies, under display faithfulness,
$\varphi(a_+) - \varphi(a_-) = \varphi(b_+) - \varphi(b_-)$. It is injective
because equal display forces balance, again by faithfulness. It is surjective
because each integer $k \geq 0$ is $\Psi([(\,\varphi^{-1}(k),\, \zd\,)])$ and
each $k < 0$ is $\Psi([(\,\zd,\, \varphi^{-1}(-k)\,)])$, where $\varphi^{-1}(k)$
is the orbit of length $k$. It preserves the operations by direct computation.
Hence $\Psi$ is a ring isomorphism.
\end{proof}

The integers also carry order, comparison, and absolute value, all native to
the signed-orbit representation.

\begin{theorem}[Order, comparison, and absolute value on $\Zd$]\label{thm:signedorder}
With distinction-only strength, $\Zd$ carries:
\begin{enumerate}[leftmargin=2em, label=\textup{(\alph*)}]
  \item a total order, $[(a_+, a_-)] \leq [(b_+, b_-)]$ if and only if
    $a_+ \oad b_- \leq_\delta b_+ \oad a_-$, well-defined on the quotient and
    compatible with addition on either side;
  \item a three-way comparison selector whose outputs are equivalent to strict
    less-than, balanced equality, and strict greater-than, computed by a
    terminating recursion;
  \item an absolute value $|[(a_+, a_-)]| = (a_+ \ominus a_-) \oad (a_- \ominus a_+)$,
    invariant under balanced equality and under negation.
\end{enumerate}
The comparison selector and absolute value are defined inside the signed-orbit
representation; the display map is used only to certify them, not to define them.
\end{theorem}

\begin{proof}
For (a), well-definedness on the quotient follows by the same cross-addition
argument as in Theorem~\ref{thm:int}, replacing equalities with the order
$\leq_\delta$ of Proposition~\ref{prop:order}, which is compatible with
addition. Totality and transitivity transfer from $\leq_\delta$. Compatibility
with integer addition follows from compatibility of $\leq_\delta$ with orbit
addition. For (b), the selector compares $a_+ \oad b_-$ with $b_+ \oad a_-$
using the terminating comparison of Proposition~\ref{prop:order} and returns one
of three outcomes; the equivalence of those outcomes with strict order and
balance is read off from the definition of $\leq_\delta$. For (c), the
expression $(a_+ \ominus a_-) \oad (a_- \ominus a_+)$ equals the larger part
minus the smaller in orbit terms, since one of the two truncated differences is
$\zd$; invariance under balance and under the swap that defines negation is a
finite check using Theorem~\ref{thm:orbit}. The transport identities
$\operatorname{cmp}(c + a, c + b) = \operatorname{cmp}(a, b)$ and
$\operatorname{cmp}(-b, -a) = \operatorname{cmp}(a, b)$ follow from the additive
compatibility and the swap symmetry.
\end{proof}

%% ===================================================================
\section{Ratio as cross-comparison: the rational field}\label{sec:ratio}

A ratio compares two integer counts. The usual construction defines a rational
as an equivalence class of pairs under cross-multiplication. We use the same
relation, expressed entirely in the integer and orbit arithmetic already built.

\begin{definition}[Ratio orbit and $\Qd$]\label{def:ratioorbit}
A \emph{ratio orbit} is a pair $(a, b)$ in which $a \in \Zd$ is the numerator
and $b \in \Nd \setminus \{\zd\}$ is a nonzero denominator length. Write
$\widehat{b} = [(b, \zd)] \in \Zd$ for the integer that lifts the denominator.
Two ratio orbits are \emph{cross-equal}, written $(a, b) \sim_\times (c, d)$,
when
\[
  a \cdot \widehat{d} \;=\; c \cdot \widehat{b}
\]
holds in $\Zd$. Set $\Qd = \{\text{ratio orbits}\} / {\sim_\times}$.
\end{definition}

For instance, the fractions $\tfrac{2}{3}$ and $\tfrac{4}{6}$ are ratio orbits
with numerators $[(2_\delta, \zd)]$ and $[(4_\delta, \zd)]$ and denominators
$3_\delta$ and $6_\delta$, where $k_\delta$ denotes the orbit of length $k$.
Cross-equality checks $2_\delta \omd 6_\delta = 4_\delta \omd 3_\delta$, that is
$12_\delta = 12_\delta$, using orbit multiplication alone. No prior notion of
fraction enters.

\begin{theorem}[The rationals form an ordered field isomorphic to $\mathbb{Q}$]\label{thm:rat}
With distinction-only strength: $\sim_\times$ is an equivalence relation; the
operations
\begin{align}
  [(a, b)] + [(c, d)] &= [(\,a \cdot \widehat{d} + c \cdot \widehat{b},\;\; b \omd d\,)],
    \label{eq:ratadd}\\
  [(a, b)] \cdot [(c, d)] &= [(\,a \cdot c,\;\; b \omd d\,)],
    \label{eq:ratmul}\\
  [(a, b)]^{-1} &= [(\,\varepsilon_\delta(a)\,\widehat{b},\;\; |a|\,)]
    \quad (a \neq 0) \label{eq:ratinv}
\end{align}
are well-defined on the quotient, where $\varepsilon_\delta(a)$ is the internal
orientation of $a$, equal to $[(\Sd(\zd), \zd)]$ when the signed-orbit
nonnegative test of Theorem~\ref{thm:signedorder} succeeds and
$[(\zd, \Sd(\zd))]$ otherwise, and $|a| \in \Nd$ is the orbit length underlying
the absolute value; under these operations $\Qd$ is an ordered field; and the
display map $\Phi \colon \Qd \to \mathbb{Q}$, $\Phi([(a, b)]) = \psi(a) / \varphi(b)$,
is an isomorphism of ordered fields.
\end{theorem}

\begin{proof}
That $\sim_\times$ is an equivalence relation follows from the commutativity and
associativity of integer multiplication, Theorem~\ref{thm:int}, together with
the fact that the lifted denominator $\widehat{b}$ is never a zero divisor
because $b \neq \zd$. Reflexivity and symmetry are immediate; transitivity uses
cancellation of the nonzero common denominator factor, which is available in
$\Zd$ since $\Zd$ is an integral domain by Theorem~\ref{thm:int}.

Well-definedness of \eqref{eq:ratadd} through \eqref{eq:ratinv} is the standard
cross-multiplication computation, carried out in $\Zd$. If $(a, b) \sim_\times (a', b')$,
then $a \cdot \widehat{b'} = a' \cdot \widehat{b}$, and substituting into each
operation and expanding with the ring laws preserves cross-equality of the
results. The reciprocal deserves a word: its denominator is the internal
absolute value $|a|$, a genuine orbit length and hence a legal denominator
because $a \neq 0$ forces $|a| \neq \zd$, and its numerator is the orientation
of $a$ times the lifted old denominator. The orientation is selected by the
internal comparison of Theorem~\ref{thm:signedorder}, not by any external sign
function, so the reciprocal is defined inside the construction. One checks
directly that $[(a, b)] \cdot [(a, b)]^{-1}$ cross-equals the unit ratio
$[(\,[(\Sd(\zd), \zd)],\, \Sd(\zd)\,)]$.

The field axioms transfer from the ring laws of $\Zd$ by computation on
representatives. The order on $\Qd$ is induced from the integer order by the rule
that $[(a, b)] \leq [(c, d)]$ when $a \cdot \widehat{d} \leq c \cdot \widehat{b}$
after normalizing denominators to be positive, which is legitimate because every
ratio orbit has a representative with positive denominator length. Finally
$\Phi$ is well-defined and injective because cross-equality matches the standard
cross-multiplication criterion for equality of fractions under display
faithfulness, surjective because every fraction $p/q$ in lowest terms is the
display of $[(\Psi^{-1}(p), \varphi^{-1}(q))]$, and operation- and
order-preserving by direct computation. Hence $\Phi$ is an isomorphism of
ordered fields.
\end{proof}

%% ===================================================================
\section{The permission surface: divisibility to greatest common divisor}\label{sec:permission}

With the number tower in place, the elementary number-theoretic surface follows
without any increase in proof budget. We collect the constructions in one
statement and then comment on what each one requires.

\begin{theorem}[Elementary permission surface]\label{thm:permission}
The following are constructed with distinction-only strength, each decidable by
a terminating procedure on finite orbit data.
\begin{enumerate}[leftmargin=2em, label=\textup{(\alph*)}]
  \item \textbf{Divisibility.} For $a, b \in \Nd$, write $a \mid b$ when there
    is $c \in \Nd$ with $b = a \omd c$. Divisibility is decidable.
  \item \textbf{Units and primality.} The only multiplicative unit of $\Nd$ is
    $\Sd(\zd)$. An orbit $p$ is \emph{prime} when $p \neq \zd$,
    $p \neq \Sd(\zd)$, and every factorization $p = a \omd b$ has
    $a = \Sd(\zd)$ or $b = \Sd(\zd)$. Primality is decidable.
  \item \textbf{Euclidean division.} For $a \in \Nd$ and $b \neq \zd$, there
    exist unique $q, r \in \Nd$ with $a = (b \omd q) \oad r$ and
    $r <_\delta b$.
  \item \textbf{Greatest common divisor.} The Euclidean algorithm on orbits
    produces $\gcd_\delta(a, b)$, the largest orbit dividing both. Two orbits
    are \emph{coprime} when $\gcd_\delta(a, b) = \Sd(\zd)$.
\end{enumerate}
Each structure has a faithful display: $\varphi$ carries divisibility,
primality, Euclidean quotient and remainder, and greatest common divisor on
$\Nd$ to the corresponding facts on $\mathbb{N}$.
\end{theorem}

\begin{proof}
For (a), the search for a witness $c$ with $b = a \omd c$ may be bounded: if
$a \neq \zd$, any witness has length at most $\varphi(b)$, so the search is
finite; if $a = \zd$, then $a \mid b$ exactly when $b = \zd$. Either way the
question is decided in finitely many steps. For (b), that $\Sd(\zd)$ is the only
unit follows because $a \omd b = \Sd(\zd)$ forces both factors to have length
one by display faithfulness; primality is then a finite search over the
candidate factors of $p$, all of which have length at most $\varphi(p)$.

For (c), existence is by induction on $a$ using repeated subtraction: if
$a <_\delta b$ then $(q, r) = (\zd, a)$; otherwise apply the hypothesis to
$a \ominus b$, which is strictly smaller, and increment the quotient. The
recursion terminates because the first argument strictly decreases at each step
and $\Nd$ is well-founded under $<_\delta$, a fact that itself follows by
induction without any appeal to a least-number principle phrased classically.
Uniqueness follows from cancellation, Theorem~\ref{thm:orbit}(e): if
$(b \omd q) \oad r = (b \omd q') \oad r'$ with $r, r' <_\delta b$, then
$q = q'$ and $r = r'$, since a difference of one in the quotient would force a
remainder of size at least $b$.

For (d), the Euclidean algorithm replaces the pair $(a, b)$ by $(b, a \bmod b)$
until the second entry is $\zd$; termination is guaranteed because the second
entry strictly decreases under $<_\delta$. The output divides both inputs and is
divisible by every common divisor, which gives the greatest-common-divisor
property. Coprimality is then the decidable condition that the output is
$\Sd(\zd)$. Display faithfulness in each case is inherited from
Theorem~\ref{thm:orbit}: $\varphi$ preserves products, hence divisibility, hence
the derived notions.
\end{proof}

It is worth pausing on what these proofs did not need. The proof of primality in
Theorem~\ref{thm:permission}(b) used orbit multiplication and the unit, and
nothing else: no set membership, no equivalence-class encoding, and no completed
infinity. The prime orbits are exactly the orbits that cannot be written as a
product of two non-unit orbits, and that property is visible in the object at
the point of construction. The same holds down the list. Order on the integers
is balance comparison. Euclidean division is repeated subtraction that
terminates because orbits are finite records. The greatest common divisor is the
ordinary Euclidean algorithm running on orbit arithmetic rather than on an
encoding of arithmetic inside a foreign primitive. In each case the structural
reason for the result is present in the construction, not recovered afterward by
decoding.

%% ===================================================================
\section{What the budget buys}\label{sec:discussion}

The constructions above are isomorphic to the familiar number systems, so the
contribution is not a new arithmetic. It is the account of what produced the
arithmetic. Three points make the account precise.

\subsection{No completed infinity}

We never posited an infinite set. Each orbit is a finite record, generated by
zero and successor, and each theorem is proved by structural induction over such
records. When we wrote $\Nd$, $\Zd$, or $\Qd$ and quantified over their elements,
we quantified over collections generated by formation rules, in the same way one
quantifies over the natural numbers themselves. A reader trained in set theory
may ask whether the inductive collection $\Nd$ is not itself an instance of the
axiom of infinity in disguise. The honest answer is that the inductive
generation of $\Nd$ is exactly as strong as accepting the natural numbers as a
generated collection, and no stronger; what we avoid is the separate
set-theoretic axiom that asserts a completed infinite set as a new object with
its own existence postulate, and the constructions never require treating $\Nd$
as a completed totality rather than as a rule for generating its members. The
distinction matters at the next layer. Completing $\Qd$ to the continuum does
require a genuinely new commitment, and that is precisely the boundary marked in
Section~\ref{sec:boundary}.

\subsection{No choice}

Every object the constructions need is produced explicitly. Representatives are
not selected from equivalence classes by a choice function; instead, all
constructions are defined on representatives and shown to respect the
equivalence, so they descend to the quotient by the principle of
Definition~\ref{def:subst}. Witnesses for divisibility, quotients and
remainders for Euclidean division, and the output of the Euclidean algorithm are
all produced by terminating procedures, not chosen from a family of candidates.
There is therefore no point at which the axiom of choice could enter.

\subsection{No excluded middle}

Every predicate the constructions rely on is decidable. Order on orbits,
equality of orbits, divisibility, primality, and coprimality are each computed
by a recursion that halts and returns a definite answer on finite data. We never
needed to assert a disjunction without being able to compute which disjunct
holds. The order trichotomy, in particular, is a computed comparison rather than
a classical case split. This is why every result carries the distinction-only
tag rather than a tag that admits classical logic.

\subsection{The display maps are instruments, not ingredients}

The maps $\varphi$, $\psi$, $\Psi$, and $\Phi$ translate the
distinction-native objects into the standard naturals, integers, and rationals,
and we proved that the relevant ones are isomorphisms. Their role is
conservativity: they certify that the construction recovers exactly the familiar
objects and produces nothing stranger. They are not used to define the native
operations. Orbit addition, signed-orbit balance, ratio cross-equality, and the
internal orientation that defines the reciprocal are all given by rules on the
native representation, and the display maps are applied afterward to verify that
those rules behave as the standard objects do. A reader who deleted the display
maps entirely would still have a complete, self-contained number tower; the maps
are there so that the reader can check the tower against the arithmetic already
known.

%% ===================================================================
\section{Where the forcing stops}\label{sec:boundary}

The constructions of this paper end at the rational field, and that endpoint is
deliberate. Everything up to and including $\Qd$ is forced by distinction and
finite repetition, with the budget held fixed at the distinction-only level. The
next natural step, completing $\Qd$ into the continuum so that limits and
suprema exist, is not forced by distinction, and the reason is structural rather
than a matter of effort. Distinction is a generative act that proceeds one step
at a time, so everything it can produce is reachable by a finite or
step-indexed generation, and the collection of all such products is countable.
The continuum is uncountable. A countable generative process cannot produce an
uncountable object, so the real line, and the completeness principle that
singles it out among ordered fields, is an addition to distinction rather than a
consequence of it.

We do not prove that claim here; it is the subject of the companion paper, which
establishes the non-forcing of the continuum with two independent arguments and
draws the consequences for the wider program. The point relevant to the present
paper is only this: the arithmetic tower is self-contained and complete at the
distinction-only level, and it is honest about its own ceiling. The boundary
between what distinction forces and what it merely permits as a further
commitment falls exactly between the rational field and its completion. Naming
that boundary precisely is part of what it means to take a single primitive
seriously, and it is the reason the arithmetic layer can be presented on its own,
as a finished result, without overstating what the primitive reaches.

%% ===================================================================
\section{Conclusion}\label{sec:conclusion}

From one act, the act of distinction, we built finite traces, a sameness
relation and its consistent companion difference relation, and a quotient
discipline. On that kernel we constructed the natural numbers as distinction
orbits, the integers as balanced pairs of orbit lengths, and the rationals as
cross-compared ratios, and we proved the complete elementary arithmetic of each:
the algebraic laws, a decidable total order, absolute value, divisibility,
primality, Euclidean division with unique quotient and remainder, greatest
common divisors, and the ordered-field axioms. Every result was established with
the same minimal proof budget, using no completed infinity, no choice, and no
classical case analysis, and we kept an explicit ledger of that budget so the
claim could be checked at each step. The integers and rationals so constructed
are isomorphic to the standard ones under faithful display maps, which we use to
certify the construction rather than to power it.

The value of the result is the audit, not a new arithmetic. It shows that the
elementary number tower is not a convention layered on top of a foundation but a
structure that the single act of distinction permits and then forces. It also
shows, by stopping where it does, exactly how far that forcing reaches: up to
the rational field and no further without a new commitment. Mathematics, at this
layer, is what distinction permits. The next paper takes up what happens when
one asks for more, and shows that the continuum is the first thing distinction
cannot supply on its own.

%% ===================================================================
\begin{thebibliography}{99}

\bibitem{Peano1889}
G.~Peano,
\emph{Arithmetices Principia, Nova Methodo Exposita},
Bocca, 1889.

\bibitem{Dedekind1888}
R.~Dedekind,
\emph{Was sind und was sollen die Zahlen?},
Vieweg, 1888.

\bibitem{Zermelo1908}
E.~Zermelo,
``Untersuchungen \"uber die Grundlagen der Mengenlehre I,''
\emph{Mathematische Annalen}, 65 (1908), 261--281.

\bibitem{MartinLof1984}
P.~Martin-L\"of,
\emph{Intuitionistic Type Theory},
Bibliopolis, 1984.

\bibitem{Bishop1967}
E.~Bishop,
\emph{Foundations of Constructive Analysis},
McGraw-Hill, 1967.

\bibitem{SpencerBrown1969}
G.~Spencer-Brown,
\emph{Laws of Form},
George Allen and Unwin, 1969.

\bibitem{Conway1976}
J.~H.~Conway,
\emph{On Numbers and Games},
Academic Press, 1976.

\bibitem{WashburnContinuum2026}
J.~Washburn,
``Where Necessity Ends: Why Distinction Forces Arithmetic but Not the
Continuum,'' draft, 2026.

\bibitem{WashburnFE2026}
J.~Washburn,
``A Functional-Equation Encoding of Logical Consistency on Continuous
Positive-Ratio Comparisons,'' draft, 2026.

\end{thebibliography}

\end{document}
